\section{Related Work}
\label{sec:related}

\paragraph{Self-Preference and Recognition.}
Prior work has documented that LLM evaluators suffer from systematic position and length biases~\citep{zheng2023judging,wang2024large}. Beyond these surface artifacts, \citet{panickssery2024llm} attribute self-preference to self-recognition, arguing that models actively detect and favor their own generations. Our findings directly refute this hypothesis across our broadest evaluation scope: all four models, all three domains, and 6{,}969 strata holding 779 distinct errors constant. Across this full grid, explicitly manipulating the authorship frame shifts verifier accuracy by less than 1\,percentage point (69.3\% neutral vs.\ 70.2\% self). Furthermore, in an 1{,}800-item self-recognition probe across all 12 model--domain cells, 11 of 12 fail to beat a conclusion-matched baseline after Bonferroni correction. In science, where apparent recognition was strongest (+22.2\,pp), controlling for prior conclusions causes the effect to collapse to an insignificant +2.8\,pp ($p = 0.73$). In code, where models cannot identify their own implementations beyond functional pass/fail outcomes, a causal style-transfer probe confirms that stripping stylistic fingerprints shifts false approval by at most $-1.5$\,pp ($p = 0.13$). Together, these results demonstrate that self-preference is driven by belief persistence---defending erroneous reasoning paths reached during generation---rather than recognition of authorship or stylistic fingerprints.

\paragraph{Evaluation Prompting and Rubrics.}
Beyond authorship, prompt design is widely treated as the primary lever for evaluator reliability. Notably, \citet{kim2023prometheus} advocate for rubric-guided evaluation, arguing that fine-grained scoring rubrics and reference criteria are essential for reliable LLM assessment. Across our full 64{,}800-verification grid, our findings reveal that the value of rubrics is strictly domain-contingent. In the code domain, where execution grounding is present, rubric prompting actively degrades verifier accuracy (overall direct scoring at 71.0\% outperforms rubrics at 68.8\%; Mistral-12B plunges from 97.7\% under direct scoring to 60--75\% under rubrics) by inducing models to overrule passing test suites with hallucinated defects. Conversely, in ungrounded reasoning domains, structured prompting remains optimal: DeepSeek-V3 achieves top accuracy in math using rubrics (69.2\%) and in science using chain-of-thought (65.8\%). Rubrics are therefore not universally beneficial: their utility depends on the domain, proving least helpful---and actively harmful---precisely where ground-truth execution signals exist.

\paragraph{Multi-Agent Deliberation, Panels, and Grounding.}
To resolve single-judge unreliability, \citet{chan2024chateval} introduce ChatEval, arguing that multi-agent debate yields superior evaluators. In our code deliberation probe ($n = 150$, DeepSeek candidates), ungrounded debate catches only 76.7\% of errors and sequential role-play catches 70.0\%, despite $28\times$ and $450\times$ higher compute costs, respectively. While narrow to this code setting, debate fails to match a single execution-grounded verifier (Qwen2.5-72B at a 93.3\% catch rate). We extend \citet{verga2024replacing}, who propose replacing individual judges with diverse panels. Across our full dataset, we reveal \emph{why} ungrounded panels fail: judges fail together on the same wrong answers rather than independently ($\kappa = 0.35\text{--}0.45$; all judges unanimously approve the same error up to $2.9\times$ more often than chance), which only decorrelate under execution grounding ($\kappa = 0.09$). Closest to our jury probe, \citet{rajan2025codexverify} introduces CodeX-Verify, reporting a 76.1\% code bug catch rate ($n = 99$) using multi-agent roles---near-identical to our ungrounded jury's 76.7\%. However, while CodeX-Verify only compares multi-agent teams against ungrounded single judges, our grounded ceiling (93.3\% catch rate) proves that the bottleneck is the absence of execution grounding, not agent count.
